% --- Archon physics formalization source begin ---
% archon:physics
% archon:covers PhyXMiniProblems/problem_phyx_mini_0230.lean
% archon:source-report reports/phyx_mini/problem_phyx_mini_0230.source.json

\chapter{Physics problem phyx\_mini\_0230}
\label{ch:PhyXMiniProblems_problem_phyx_mini_0230}

\paragraph{Problem source.}
\#\# Physical scenario

An object attached to a spring vibrates with simple harmonic motion as described by Figure.

\#\# Question

For this motion, find the maximum acceleration.

\#\# Answer choices

- A: A: \$4.39\textbackslash{}

- B: B: \$5.93\textbackslash{}

- C: C: \$3.93\textbackslash{}

- D: D: \$4.93\textbackslash{}

\#\# Figure caption (auxiliary; use the image as primary evidence)

The image displays a graph featuring a sinusoidal wave. The horizontal axis is labeled \textbackslash{}( t \textbackslash{}) (s), representing time in seconds, and the vertical axis is labeled \textbackslash{}( x \textbackslash{}) (cm), representing displacement in centimeters. The wave oscillates between \textbackslash{}( 2.00 \textbackslash{}) and \textbackslash{}(-2.00\textbackslash{}) on the vertical axis.

Key points along the wave are marked with black dots at the following coordinates: (0, 0), (1, 2), (3, -2), (5, 2), and (6, 0). The wave completes one full cycle from \textbackslash{}( t = 0 \textbackslash{}) to \textbackslash{}( t = 6 \textbackslash{}).

The wave is visually delineated with a brown curve and the gridlines are visible, aiding with precision in reading the graph's marked points and structure.

\#\# Dataset metadata

Wave Properties; Physical Model Grounding Reasoning, Numerical Reasoning

\paragraph{Recorded answer/context.}
D: D: \$4.93\textbackslash{}

\paragraph{Figure/image path.}
/root/proposal\_for\_physic/science-mango/phyx\_mini\_run/phyx\_data/test\_image/230.png

\paragraph{Formalization target.}
create a compiling Lean file with sorry bodies at `PhyXMiniProblems/problem\_phyx\_mini\_0230.lean`. The Lean declarations must preserve the physical quantities, dimensions or dimensional roles, figure labels, governing-law hypotheses, and final relation expressed by this problem.
Use Mathlib/Physlib names found through LeanExplore where available. If a physics API is missing, introduce faithful local abstractions rather than scalar placeholder aliases.

\begin{theorem}[Physics formalization target]
\label{thm:physics:phyx_mini_0230:target}
\lean{PhyXMiniProblems.ProblemPhyXMini0230.problem_phyx_mini_0230}
\uses{def:physics:phyx-mini-0230:phyxminiproblems-problemphyxmini0230-accelerationquantity, def:physics:phyx-mini-0230:phyxminiproblems-problemphyxmini0230-lengthreadout, def:physics:phyx-mini-0230:phyxminiproblems-problemphyxmini0230-timereadout, def:physics:phyx-mini-0230:phyxminiproblems-problemphyxmini0230-angularfrequencyreadout, def:physics:phyx-mini-0230:phyxminiproblems-problemphyxmini0230-accelerationreadout, def:physics:phyx-mini-0230:phyxminiproblems-problemphyxmini0230-springoscillationsetup, def:physics:phyx-mini-0230:phyxminiproblems-problemphyxmini0230-matchesspringmotionfigure, def:physics:phyx-mini-0230:phyxminiproblems-problemphyxmini0230-hasphysicalspringmotionparameters, def:physics:phyx-mini-0230:phyxminiproblems-problemphyxmini0230-satisfiessimpleharmonicmotionlaws, def:physics:phyx-mini-0230:phyxminiproblems-problemphyxmini0230-ismaximumaccelerationmagnitude, def:physics:phyx-mini-0230:phyxminiproblems-problemphyxmini0230-answerchoice, def:physics:phyx-mini-0230:phyxminiproblems-problemphyxmini0230-answerchoiceincentimeterspersecondsquared}
The assigned autoformalize agent should translate this physics problem into Lean declarations in the covered file, with theorem and lemma proof bodies written as `by sorry`.
\end{theorem}
\begin{proof}
This is an autoformalization task, not a proof task. Produce statements that can later be proved by the physics prover without weakening the physical model.
\end{proof}

% --- Archon declaration topology begin ---
\section{Declaration topology}
\begin{definition}[Length Quantity]
  \label{def:physics:phyx-mini-0230:phyxminiproblems-problemphyxmini0230-lengthquantity}
  \lean{PhyXMiniProblems.ProblemPhyXMini0230.LengthQuantity}
  A signed displacement or oscillation amplitude
\end{definition}
\begin{proof}
  This is the declared model definition of Length Quantity.
\end{proof}
\begin{definition}[Time Quantity]
  \label{def:physics:phyx-mini-0230:phyxminiproblems-problemphyxmini0230-timequantity}
  \lean{PhyXMiniProblems.ProblemPhyXMini0230.TimeQuantity}
  A physical time or period
\end{definition}
\begin{proof}
  This is the declared model definition of Time Quantity.
\end{proof}
\begin{definition}[Angular Frequency Quantity]
  \label{def:physics:phyx-mini-0230:phyxminiproblems-problemphyxmini0230-angularfrequencyquantity}
  \lean{PhyXMiniProblems.ProblemPhyXMini0230.AngularFrequencyQuantity}
  Angular frequency; radians are dimensionless, so its dimension is inverse time
\end{definition}
\begin{proof}
  This is the declared model definition of Angular Frequency Quantity.
\end{proof}
\begin{definition}[Acceleration Quantity]
  \label{def:physics:phyx-mini-0230:phyxminiproblems-problemphyxmini0230-accelerationquantity}
  \lean{PhyXMiniProblems.ProblemPhyXMini0230.AccelerationQuantity}
  A signed one-dimensional acceleration
\end{definition}
\begin{proof}
  This is the declared model definition of Acceleration Quantity.
\end{proof}
\begin{definition}[length Readout]
  \label{def:physics:phyx-mini-0230:phyxminiproblems-problemphyxmini0230-lengthreadout}
  \lean{PhyXMiniProblems.ProblemPhyXMini0230.lengthReadout}
  \uses{def:physics:phyx-mini-0230:phyxminiproblems-problemphyxmini0230-lengthquantity}
  Read a physical length in a selected length unit
\end{definition}
\begin{proof}
  This is the declared model definition of length Readout.
\end{proof}
\begin{definition}[time Readout]
  \label{def:physics:phyx-mini-0230:phyxminiproblems-problemphyxmini0230-timereadout}
  \lean{PhyXMiniProblems.ProblemPhyXMini0230.timeReadout}
  \uses{def:physics:phyx-mini-0230:phyxminiproblems-problemphyxmini0230-timequantity}
  Read a physical time in a selected time unit
\end{definition}
\begin{proof}
  This is the declared model definition of time Readout.
\end{proof}
\begin{definition}[angular Frequency Readout]
  \label{def:physics:phyx-mini-0230:phyxminiproblems-problemphyxmini0230-angularfrequencyreadout}
  \lean{PhyXMiniProblems.ProblemPhyXMini0230.angularFrequencyReadout}
  \uses{def:physics:phyx-mini-0230:phyxminiproblems-problemphyxmini0230-angularfrequencyquantity}
  Read angular frequency in radians per selected time unit
\end{definition}
\begin{proof}
  This is the declared model definition of angular Frequency Readout.
\end{proof}
\begin{definition}[acceleration Readout]
  \label{def:physics:phyx-mini-0230:phyxminiproblems-problemphyxmini0230-accelerationreadout}
  \lean{PhyXMiniProblems.ProblemPhyXMini0230.accelerationReadout}
  \uses{def:physics:phyx-mini-0230:phyxminiproblems-problemphyxmini0230-accelerationquantity}
  Read acceleration in a coherent selected length/time unit system
\end{definition}
\begin{proof}
  This is the declared model definition of acceleration Readout.
\end{proof}
\begin{definition}[Graph Axis]
  \label{def:physics:phyx-mini-0230:phyxminiproblems-problemphyxmini0230-graphaxis}
  \lean{PhyXMiniProblems.ProblemPhyXMini0230.GraphAxis}
  The two coordinate axes shown in the graph
\end{definition}
\begin{proof}
  This is the declared model definition of Graph Axis.
\end{proof}
\begin{definition}[Axis Quantity]
  \label{def:physics:phyx-mini-0230:phyxminiproblems-problemphyxmini0230-axisquantity}
  \lean{PhyXMiniProblems.ProblemPhyXMini0230.AxisQuantity}
  The physical quantity printed on each graph axis
\end{definition}
\begin{proof}
  This is the declared model definition of Axis Quantity.
\end{proof}
\begin{definition}[Figure Sample]
  \label{def:physics:phyx-mini-0230:phyxminiproblems-problemphyxmini0230-figuresample}
  \lean{PhyXMiniProblems.ProblemPhyXMini0230.FigureSample}
  The seven black sample points explicitly marked in the primary image
\end{definition}
\begin{proof}
  This is the declared model definition of Figure Sample.
\end{proof}
\begin{definition}[figure Time In Seconds]
  \label{def:physics:phyx-mini-0230:phyxminiproblems-problemphyxmini0230-figuretimeinseconds}
  \lean{PhyXMiniProblems.ProblemPhyXMini0230.figureTimeInSeconds}
  \uses{def:physics:phyx-mini-0230:phyxminiproblems-problemphyxmini0230-figuresample}
  Printed time coordinate, in seconds, of each marked point
\end{definition}
\begin{proof}
  This is the declared model definition of figure Time In Seconds.
\end{proof}
\begin{definition}[figure Displacement In Centimeters]
  \label{def:physics:phyx-mini-0230:phyxminiproblems-problemphyxmini0230-figuredisplacementincentimeters}
  \lean{PhyXMiniProblems.ProblemPhyXMini0230.figureDisplacementInCentimeters}
  \uses{def:physics:phyx-mini-0230:phyxminiproblems-problemphyxmini0230-figuresample}
  Printed displacement coordinate, in centimeters, of each marked point
\end{definition}
\begin{proof}
  This is the declared model definition of figure Displacement In Centimeters.
\end{proof}
\begin{definition}[Restoring Element]
  \label{def:physics:phyx-mini-0230:phyxminiproblems-problemphyxmini0230-restoringelement}
  \lean{PhyXMiniProblems.ProblemPhyXMini0230.RestoringElement}
  The qualitative restoring element named in the problem statement
\end{definition}
\begin{proof}
  This is the declared model definition of Restoring Element.
\end{proof}
\begin{definition}[Motion Regime]
  \label{def:physics:phyx-mini-0230:phyxminiproblems-problemphyxmini0230-motionregime}
  \lean{PhyXMiniProblems.ProblemPhyXMini0230.MotionRegime}
  The motion regime asserted by the problem statement
\end{definition}
\begin{proof}
  This is the declared model definition of Motion Regime.
\end{proof}
\begin{definition}[Displacement Time Graph]
  \label{def:physics:phyx-mini-0230:phyxminiproblems-problemphyxmini0230-displacementtimegraph}
  \lean{PhyXMiniProblems.ProblemPhyXMini0230.DisplacementTimeGraph}
  \uses{def:physics:phyx-mini-0230:phyxminiproblems-problemphyxmini0230-timequantity, def:physics:phyx-mini-0230:phyxminiproblems-problemphyxmini0230-graphaxis, def:physics:phyx-mini-0230:phyxminiproblems-problemphyxmini0230-axisquantity, def:physics:phyx-mini-0230:phyxminiproblems-problemphyxmini0230-figuresample}
  Axis labels, units, and qualitative rendering features of the graph
\end{definition}
\begin{proof}
  This is the declared model definition of Displacement Time Graph.
\end{proof}
\begin{definition}[Spring Oscillation Setup]
  \label{def:physics:phyx-mini-0230:phyxminiproblems-problemphyxmini0230-springoscillationsetup}
  \lean{PhyXMiniProblems.ProblemPhyXMini0230.SpringOscillationSetup}
  \uses{def:physics:phyx-mini-0230:phyxminiproblems-problemphyxmini0230-lengthquantity, def:physics:phyx-mini-0230:phyxminiproblems-problemphyxmini0230-timequantity, def:physics:phyx-mini-0230:phyxminiproblems-problemphyxmini0230-angularfrequencyquantity, def:physics:phyx-mini-0230:phyxminiproblems-problemphyxmini0230-accelerationquantity, def:physics:phyx-mini-0230:phyxminiproblems-problemphyxmini0230-restoringelement, def:physics:phyx-mini-0230:phyxminiproblems-problemphyxmini0230-motionregime, def:physics:phyx-mini-0230:phyxminiproblems-problemphyxmini0230-displacementtimegraph}
  The spring oscillator and its dimensionful kinematic quantities. The acceleration function is not defined from the requested answer. It is an independent physical quantity linked to displacement only by the governing simple-harmonic-motion laws below
\end{definition}
\begin{proof}
  This is the declared model definition of Spring Oscillation Setup.
\end{proof}
\begin{definition}[Matches Spring Motion Figure]
  \label{def:physics:phyx-mini-0230:phyxminiproblems-problemphyxmini0230-matchesspringmotionfigure}
  \lean{PhyXMiniProblems.ProblemPhyXMini0230.MatchesSpringMotionFigure}
  \uses{def:physics:phyx-mini-0230:phyxminiproblems-problemphyxmini0230-lengthreadout, def:physics:phyx-mini-0230:phyxminiproblems-problemphyxmini0230-timereadout, def:physics:phyx-mini-0230:phyxminiproblems-problemphyxmini0230-figuresample, def:physics:phyx-mini-0230:phyxminiproblems-problemphyxmini0230-figuretimeinseconds, def:physics:phyx-mini-0230:phyxminiproblems-problemphyxmini0230-figuredisplacementincentimeters, def:physics:phyx-mini-0230:phyxminiproblems-problemphyxmini0230-springoscillationsetup}
  Scenario and primary-image evidence. The amplitude and period here are graph readouts, not conclusions about acceleration. The period is the separation of the consecutive positive peaks at samples t1 and t5
\end{definition}
\begin{proof}
  This is the declared model definition of Matches Spring Motion Figure.
\end{proof}
\begin{definition}[Has Physical Spring Motion Parameters]
  \label{def:physics:phyx-mini-0230:phyxminiproblems-problemphyxmini0230-hasphysicalspringmotionparameters}
  \lean{PhyXMiniProblems.ProblemPhyXMini0230.HasPhysicalSpringMotionParameters}
  \uses{def:physics:phyx-mini-0230:phyxminiproblems-problemphyxmini0230-lengthreadout, def:physics:phyx-mini-0230:phyxminiproblems-problemphyxmini0230-timereadout, def:physics:phyx-mini-0230:phyxminiproblems-problemphyxmini0230-angularfrequencyreadout, def:physics:phyx-mini-0230:phyxminiproblems-problemphyxmini0230-springoscillationsetup}
  Positivity conditions for the nondegenerate oscillator shown in the graph
\end{definition}
\begin{proof}
  This is the declared model definition of Has Physical Spring Motion Parameters.
\end{proof}
\begin{definition}[Satisfies Simple Harmonic Motion Laws]
  \label{def:physics:phyx-mini-0230:phyxminiproblems-problemphyxmini0230-satisfiessimpleharmonicmotionlaws}
  \lean{PhyXMiniProblems.ProblemPhyXMini0230.SatisfiesSimpleHarmonicMotionLaws}
  \uses{def:physics:phyx-mini-0230:phyxminiproblems-problemphyxmini0230-timequantity, def:physics:phyx-mini-0230:phyxminiproblems-problemphyxmini0230-lengthreadout, def:physics:phyx-mini-0230:phyxminiproblems-problemphyxmini0230-timereadout, def:physics:phyx-mini-0230:phyxminiproblems-problemphyxmini0230-angularfrequencyreadout, def:physics:phyx-mini-0230:phyxminiproblems-problemphyxmini0230-accelerationreadout, def:physics:phyx-mini-0230:phyxminiproblems-problemphyxmini0230-springoscillationsetup}
  Governing laws for the chosen phase of the simple harmonic motion. The graph crosses equilibrium at t = 0 while moving toward its positive peak, hence the sine phase. The other fields state ω = 2π/T and a = -ω²x in the coherent centimeter/second readouts. No maximum acceleration or answer-choice value occurs in this premise structure
\end{definition}
\begin{proof}
  This is the declared model definition of Satisfies Simple Harmonic Motion Laws.
\end{proof}
\begin{definition}[Is Maximum Acceleration Magnitude]
  \label{def:physics:phyx-mini-0230:phyxminiproblems-problemphyxmini0230-ismaximumaccelerationmagnitude}
  \lean{PhyXMiniProblems.ProblemPhyXMini0230.IsMaximumAccelerationMagnitude}
  \uses{def:physics:phyx-mini-0230:phyxminiproblems-problemphyxmini0230-timequantity, def:physics:phyx-mini-0230:phyxminiproblems-problemphyxmini0230-accelerationquantity, def:physics:phyx-mini-0230:phyxminiproblems-problemphyxmini0230-accelerationreadout, def:physics:phyx-mini-0230:phyxminiproblems-problemphyxmini0230-springoscillationsetup}
  A physical acceleration magnitude that is attained and bounds the motion
\end{definition}
\begin{proof}
  This is the declared model definition of Is Maximum Acceleration Magnitude.
\end{proof}
\begin{definition}[Answer Choice]
  \label{def:physics:phyx-mini-0230:phyxminiproblems-problemphyxmini0230-answerchoice}
  \lean{PhyXMiniProblems.ProblemPhyXMini0230.AnswerChoice}
  Labels of the four displayed answer choices
\end{definition}
\begin{proof}
  This is the declared model definition of Answer Choice.
\end{proof}
\begin{definition}[answer Choice In Centimeters Per Second Squared]
  \label{def:physics:phyx-mini-0230:phyxminiproblems-problemphyxmini0230-answerchoiceincentimeterspersecondsquared}
  \lean{PhyXMiniProblems.ProblemPhyXMini0230.answerChoiceInCentimetersPerSecondSquared}
  \uses{def:physics:phyx-mini-0230:phyxminiproblems-problemphyxmini0230-answerchoice}
  Displayed answer readouts, inferred to be in centimeters per second squared
\end{definition}
\begin{proof}
  This is the declared model definition of answer Choice In Centimeters Per Second Squared.
\end{proof}
% --- Archon declaration topology end ---
% --- Archon physics formalization source end ---
